\documentclass[11pt]{article}
% Macros and theorem environments: \newcommand with arguments and an
% optional default, \def, \newenvironment, \DeclareMathOperator, and the
% \newtheorem family with proofs. A macro a formula uses is prepended to
% that formula's source, so the math engine expands it exactly as the text
% does — the definition is written once and holds in both modes.
\usepackage{amsmath}
\usepackage{amsthm}
\usepackage{amssymb}
\newcommand{\R}{\mathbb{R}}
\newcommand{\N}{\mathbb{N}}
\newcommand{\norm}[1]{\left\lVert #1 \right\rVert}
\newcommand{\inner}[2]{\langle #1, #2 \rangle}
\newcommand{\greet}[2][Dear]{#1 #2,}
\def\framework{UltraCanvas}
\DeclareMathOperator{\tr}{tr}
\DeclareMathOperator*{\argmin}{arg\,min}
\newenvironment{aside}{\begin{quote}\itshape}{\end{quote}}
\newtheorem{theorem}{Theorem}
\newtheorem{lemma}[theorem]{Lemma}
\newtheorem{definition}{Definition}
\title{Macros, Operators and Theorems}
\author{The \framework{} Team}
\date{September 2026}
\begin{document}
\maketitle

\begin{abstract}
Every definition in the preamble above is expanded by the importer, in text
and in mathematics alike. \framework{} is a \verb|\def|; \texttt{Dear
reader,} below is a \verb|\newcommand| with an optional argument; and
\texttt{tr} and \texttt{arg\,min} are operators declared with
\verb|\DeclareMathOperator|, which the math engine sets upright and spaces
as operators rather than as a product of letters.
\end{abstract}

\section{Text macros}
\greet{reader} a macro with an optional argument uses its default unless one
is given: \greet[Hello]{reader} shows the other branch. A macro may be used
inside another, and a \verb|\def| such as \verb|\framework| takes no
arguments at all --- this paragraph mentions \framework{} twice.

\begin{aside}
\texttt{aside} is a \verb|\newenvironment| whose body is wrapped in a quote
and set in italics. The importer expands the begin and end texts around the
body, so the environment needs no support of its own.
\end{aside}

\section{Macros in mathematics}\label{sec:math}
The same \verb|\norm| and \verb|\inner| macros are used in the formulas
below; nothing declares them a second time.
\begin{equation}\label{eq:cauchy}
  \left| \inner{u}{v} \right| \le \norm{u} \, \norm{v},
  \qquad u, v \in \R^n
\end{equation}
Equation~\eqref{eq:cauchy} is the Cauchy--Schwarz inequality. Operators
declared in the preamble keep their spacing in display style,
\[
  \tr(AB) = \tr(BA), \qquad
  \hat{\theta} = \argmin_{\theta \in \R^p} \norm{y - X\theta}^2 ,
\]
and the starred \verb|\DeclareMathOperator*| is what puts the subscript
under \texttt{arg\,min} rather than beside it.

\section{Theorems}
\texttt{\textbackslash newtheorem} declares an environment, its printed
name and --- optionally --- the counter it shares. Here \texttt{lemma}
shares the theorem counter while \texttt{definition} counts on its own.

\begin{definition}[Norm]\label{def:norm}
A norm on $\R^n$ is a map $\norm{\cdot} : \R^n \to \R$ that is positive
definite, absolutely homogeneous and subadditive.
\end{definition}

\begin{theorem}[Cauchy--Schwarz]\label{thm:cs}
For all $u, v \in \R^n$, $\left| \inner{u}{v} \right| \le \norm{u} \norm{v}$,
with equality exactly when $u$ and $v$ are linearly dependent.
\end{theorem}

\begin{proof}
For $t \in \R$ expand $0 \le \norm{u - t v}^2 = \norm{u}^2 - 2t \inner{u}{v}
+ t^2 \norm{v}^2$. The right-hand side is a quadratic in $t$ that is never
negative, so its discriminant is not positive:
$4 \inner{u}{v}^2 - 4 \norm{u}^2 \norm{v}^2 \le 0$.
\end{proof}

\begin{lemma}\label{lem:triangle}
The triangle inequality $\norm{u + v} \le \norm{u} + \norm{v}$ follows from
Theorem~\ref{thm:cs}, and this lemma is numbered~\ref{lem:triangle} because
it shares the theorem counter.
\end{lemma}

Definition~\ref{def:norm} counts separately, which is why it is number~1
while the theorem above it is also number~1.
\end{document}
